\documentclass[letterpaper]{article} % DO NOT CHANGE THIS
\pdfinfoomitdate=1
\pdftrailerid{}
\pdfsuppressptexinfo=-1
\pdfinfo{/Creator () /Producer ()}
\usepackage[submission]{aaai2027}  % DO NOT CHANGE THIS
\usepackage[hyphens]{url}  % DO NOT CHANGE THIS
\usepackage{graphicx} % DO NOT CHANGE THIS
\urlstyle{rm} % DO NOT CHANGE THIS
\def\UrlFont{\rm}  % DO NOT CHANGE THIS
\usepackage{natbib}  % DO NOT CHANGE THIS AND DO NOT ADD ANY OPTIONS TO IT
\usepackage{caption} % DO NOT CHANGE THIS AND DO NOT ADD ANY OPTIONS TO IT
\frenchspacing  % DO NOT CHANGE THIS

\usepackage{amsmath}
\usepackage{booktabs}
\usepackage{algorithm}
\usepackage{algorithmic}

\pdfinfo{
/TemplateVersion (2027.1)
}

\setcounter{secnumdepth}{1}
\numberwithin{equation}{section}
\numberwithin{table}{section}

\title{Supplementary Material}
\affiliations{}

\begin{document}

\maketitle
\appendix

This supplement provides the technical details referenced by the main paper. All symbols follow the definitions in the main paper.

\section{Method-Aware Latency}
\label{app:latency-prior}

\subsection{Scope and notation}

For a Linear module, we use the convention \(Y=XW^\top\), where \(X\),
\(W\), and \(Y\) have shapes \(m\times k\), \(n\times k\), and \(m\times n\),
respectively. The latency predictor targets the complete online module
forward, including representation-specific activation processing, weight and
scale reads, sparse metadata access and output writes. Offline
checkpoint conversion and weight packing are not included because they are
performed before serving.

For compression method \(c\), \(F_c(m,n,k)\) is the effective floating-point
work and \(D_c(m,n,k)\) is the logical data movement used by the Roofline
prior. These quantities are not intended to reproduce every physical memory
transaction.

\subsection{Method-specific work and data movement}

\paragraph{Dense BF16.}
The dense kernel reads BF16 activations and weights and writes a BF16 output:
\begin{equation}
\begin{aligned}
F_{\mathrm{dbf16}}(m,n,k) &= 2mnk,\\
D_{\mathrm{dbf16}}(m,n,k) &= 2mk+2nk+2mn.
\end{aligned}
\end{equation}

\paragraph{Sparse BF16.}
The 2:4 kernel retains two BF16 values in every group of four weights, halving
the effective multiply--accumulate count. A dense group occupies \(64\) bits,
whereas its sparse representation contains \(32\) value bits and \(4\)
metadata bits. The packed weight traffic is therefore
\begin{equation}
2nk\big/(64/36)=\frac{9}{8}nk,
\end{equation}
which gives
\begin{equation}
\begin{aligned}
F_{\mathrm{sbf16}}(m,n,k) &= mnk,\\
D_{\mathrm{sbf16}}(m,n,k) &= 2mk+\frac{9}{8}nk+2mn.
\end{aligned}
\end{equation}


\paragraph{Dense NVFP4.}
The dense W4A4 path reads the BF16 input, produces a packed FP4 activation,
reads FP4 weights and their associated scales, and writes a BF16 output:
\begin{equation}
\begin{aligned}
F_{\mathrm{dnvfp4}}(m,n,k) &= 2mnk,\\
D_{\mathrm{dnvfp4}}(m,n,k)
&=2mk+0.5mk+0.5nk\\
&\quad+\frac{mk+nk}{16}+2mn.
\end{aligned}
\end{equation}
The terms \(0.5mk\) and \(0.5nk\) represent packed 4-bit activation and
weight values. NVFP4 uses one FP8 scale for each group of \(16\) values and
one FP32 global scale for the entire tensor. Accordingly,
\((mk+nk)/16\) accounts for the FP8 scale traffic of the activations and
weights. We omit the single FP32 global scale per tensor because its constant
four-byte cost is negligible relative to the matrix data.

\paragraph{Sparse NVFP4.}
Due to underlying hardware constraints and the requirements of the CUTLASS
sparse kernel, Sparse NVFP4 uses one FP8 scale for each group of \(32\)
values rather than the \(16\)-value groups used by dense NVFP4. It applies
50\% semi-structured sparsity using the kernel's pairwise 4:8 layout: four FP4
values are retained in each group of eight weights. Consequently, the
effective work is halved and the packed weight values occupy \(0.25nk\)
bytes. The sparse layout additionally uses \(4\) position bits per group,
equivalent to \(nk/16\) bytes. Including the FP8 activation and weight scales,
while again omitting the tensor-level FP32 global scales, gives
\begin{equation}
\begin{aligned}
F_{\mathrm{snvfp4}}(m,n,k) &= mnk,\\
D_{\mathrm{snvfp4}}(m,n,k)
&=2mk+0.5mk+\frac{mk}{32}\\
&\quad+0.25nk+\frac{nk}{16}+\frac{nk}{32}+2mn\\
&=2.5mk+\frac{mk}{32}+\frac{11}{32}nk+2mn.
\end{aligned}
\end{equation}
This description distinguishes the logical 50\% sparsity used for comparison
from the physical 4:8 grouping required by the sparse FP4 kernel.

\paragraph{W4A16 Marlin.}
The weight-only path retains BF16 activations and outputs while reading packed
4-bit weights and their scales:
\begin{equation}
\begin{aligned}
F_{\mathrm{marlin}}(m,n,k) &= 2mnk,\\
D_{\mathrm{marlin}}(m,n,k)
&=2mk+0.5nk+\frac{nk}{16}+2mn.
\end{aligned}
\end{equation}

\section{Residual Correction}
\label{app:latency-residual}

\subsection{Residual target and global features}

For a profiled shape \(x_j=(m_j,n_j,k_j)\), the logarithmic residual is
\begin{equation}
\delta_{c,j}
=
\log\frac{\ell^{\mathrm{meas}}_{c,j}}
{\ell_c^{\mathrm{base}}(x_j)}.
\end{equation}
The global predictor is a ridge regressor
\(\widehat{\delta}^{\mathrm{global}}_c(x)=\theta_c^\top\phi_c(x)\).
Its feature vector includes log dimensions and aspect ratios, kernel-specific
tile counts and tail-utilization statistics, and indicators for small or
highly skewed matrix shapes.
Tile and tail features use the native tile geometry of each kernel family
rather than a shared tile size across all compression methods.

\subsection{Local interpolation}

The local predictor operates in log-shape space. Let
\(u(x)=(\log_2m,\log_2n,\log_2k)\), where a unit change in any coordinate
corresponds to doubling the associated matrix dimension. For kernel family
\(c\), we use the Euclidean distance
\(d_c(x,x_j)=\|u(x)-u(x_j)\|_2\). Let
\(\mathcal{N}_{12,c}(x)\) be the 12 nearest profiled shapes. The local
residual estimate is
\begin{equation}
\widehat{\delta}^{\mathrm{local}}_c(x)
=
\frac{
\displaystyle\sum_{j\in\mathcal{N}_{12,c}(x)}
\frac{\delta_{c,j}}{d_c(x,x_j)+10^{-6}}
}{
\displaystyle\sum_{j\in\mathcal{N}_{12,c}(x)}
\frac{1}{d_c(x,x_j)+10^{-6}}
}.
\end{equation}

\subsection{Adaptive mixing}

Let \(d_{\min,c}(x)\) be the distance to the nearest profiled shape. We assign
more weight to local interpolation near observed shapes and gradually shift
toward the regularized global predictor when extrapolating:
\begin{equation}
\lambda_c(x)
=
\operatorname{clip}\!\left(
\frac{2.5}{2.5+d_{\min,c}(x)},\,0.25,\,0.95
\right),
\end{equation}
\begin{equation}
\widehat{\delta}_c(x)
=
\lambda_c(x)\widehat{\delta}^{\mathrm{local}}_c(x)
+\bigl[1-\lambda_c(x)\bigr]
\widehat{\delta}^{\mathrm{global}}_c(x).
\end{equation}
The calibrated module latency is
\begin{equation}
\widehat{\ell}_c(x)
=
\ell_c^{\mathrm{base}}(x)\exp\!\bigl(\widehat{\delta}_c(x)\bigr).
\end{equation}
When an exact profiled shape is available, the implementation returns its
measured median directly. The learned correction is used for shapes absent
from the profile table.

% \section{Structured Quality-Cost Construction}
% \label{app:quality-cost}

% \subsection{Shared structural factors}

% The policy remains configurable at individual Linear modules; structural
% labels only determine which sensitivity parameters are shared. For both
% evaluated 32-layer Llama models, we use four contiguous depth regions:
% \begin{equation}
% b_i=\left\lfloor\frac{l_i}{8}\right\rfloor
% \in\{0,1,2,3\},
% \end{equation}
% corresponding to layers \(0\)--\(7\), \(8\)--\(15\), \(16\)--\(23\), and
% \(24\)--\(31\). Here \(l_i\) is the layer index of module \(v_i\).
% The module-type label \(t_i\) takes one of four values:
% \texttt{qkv\_proj}, \texttt{o\_proj}, \texttt{gate\_up\_proj}, and
% \texttt{down\_proj}. Modules within the same region or type share the
% corresponding sensitivity term, but the optimizer still selects their
% compression actions independently.

% \subsection{Policy-level feature aggregation}

% For policy \(\pi\), define the selected local-perturbation mass associated
% with phase \(r\), method \(c\), depth region \(b\), and module type \(t\) as
% \begin{equation}
% X_{r,c,b,t}(\pi)
% =
% \sum_{i=1}^{M}
% \mathbf{1}[b_i=b]\,
% \mathbf{1}[t_i=t]\,
% \mathbf{1}[a_i^r=c]\,
% e_{r,i,c}.
% \end{equation}
% The structured model in the main paper can then be written equivalently as
% \begin{equation}
% \begin{aligned}
% \widehat{\Delta\mathrm{NLL}}(\pi)
% &=q_0+\frac{1}{s_{\mathrm{feat}}}
% \sum_{r,c,b,t}X_{r,c,b,t}(\pi)\\
% &\quad{}\times
% \operatorname{ReLU}
% \!\left(g+\rho_r+u_c+\beta_b+\tau_t\right).
% \end{aligned}
% \end{equation}
% The normalization constant is the average selected feature mass over the
% calibration policies \(\Pi_{\mathrm{cal}}\):
% \begin{equation}
% s_{\mathrm{feat}}
% =
% \frac{1}{|\Pi_{\mathrm{cal}}|}
% \sum_{\pi\in\Pi_{\mathrm{cal}}}
% \sum_{r,c,b,t}X_{r,c,b,t}(\pi).
% \end{equation}
% The prefill-only and prefill--decode scenarios are fitted separately because
% they use different phase sets and policy-level NLL labels. In both cases, the
% parameters are learned from complete heterogeneous policies by minimizing
% regularized mean-squared prediction error:
% \begin{equation}
% \frac{1}{|\Pi_{\mathrm{cal}}|}
% \sum_{\pi\in\Pi_{\mathrm{cal}}}
% \left[
% \widehat{\Delta\mathrm{NLL}}(\pi)
% -\Delta\mathrm{NLL}(\pi)
% \right]^2
% +\gamma\|\Theta\|_2^2.
% \end{equation}
% The intercept \(q_0\) is left unregularized. The resulting
% \(q_{r,i,c}\) is an additive optimization surrogate calibrated from joint
% policies; it should not be interpreted as an independently measured causal
% NLL increase from compressing only module \(v_i\).

\section{Budgeted Policy Optimization}
\label{app:optimization}

For module \(v_i\), a legal action is \(a=(a^p,a^d)\in\mathcal{A}_i\) in a
prefill--decode scenario, or a single prefill action in a prefill-only
scenario. Its latency and variable quality costs are
\begin{equation}
\ell_i(a;s)
=
\widehat{\ell}_{a^p}(BS,n_i,k_i)
+(T-1)\widehat{\ell}_{a^d}(B,n_i,k_i),
\end{equation}
\begin{equation}
q_i(a;s)
=
\sum_{r\in\mathcal{R}}q_{r,i,a^r},
\end{equation}
with the decoding term omitted for prefill-only inference. The
policy-independent intercept \(q_0\) is subtracted once from the quality
budget before optimization.

Let \(\Delta_q\) be the quality discretization step,
\(\bar q_i(a)=\lceil q_i(a;s)/\Delta_q\rceil\), and \(H\) the discretized
budget. Define \(D_i(h)\) as the minimum latency attainable by the first
\(i\) modules with cumulative quality cost at most \(h\). The recurrence is
\begin{equation}
D_i(h)
=
\min_{\substack{a\in\mathcal{A}_i\\\bar q_i(a)\leq h}}
\left\{
D_{i-1}\!\left(h-\bar q_i(a)\right)+\ell_i(a;s)
\right\},
\end{equation}
with \(D_0(h)=0\) for \(h=0,\ldots,H\). Storing the minimizing action for
each state permits policy reconstruction. The solver requires
\(O\!\left(H\sum_i|\mathcal{A}_i|\right)\) time and \(O(H)\) working memory.
Solving across quality budgets and removing dominated or duplicate policies
produces the predicted frontier.

\begin{algorithm}[t]
\caption{Quality-budgeted policy optimization}
\label{alg:budgeted-policy-optimization}
\small
\begin{algorithmic}[1]
\STATE \textbf{Input:} action sets \(\{\mathcal{A}_i\}_{i=1}^{M}\);
latency costs \(\ell_i(a;s)\); quality costs \(q_i(a;s)\);
budget \(\epsilon\); intercept \(q_0\); step \(\Delta_q\)
\STATE \textbf{Output:} policy \(\pi=(a_1,\ldots,a_M)\)
\STATE \(H\gets
\left\lfloor\max(0,\epsilon-q_0)/\Delta_q\right\rfloor\)
\STATE \(D[h]\gets0\) for \(h=0,\ldots,H\)
\FOR{\(i=1\) \TO \(M\)}
    \STATE \(D_{\mathrm{new}}[h]\gets\infty\) for \(h=0,\ldots,H\)
    \FOR{\(h=0\) \TO \(H\)}
        \FORALL{\(a\in\mathcal{A}_i\)}
            \STATE \(b\gets\left\lceil q_i(a;s)/\Delta_q\right\rceil\)
            \IF{\(b\leq h\) \AND
                \(D[h-b]+\ell_i(a;s)<D_{\mathrm{new}}[h]\)}
                \STATE \(D_{\mathrm{new}}[h]\gets
                    D[h-b]+\ell_i(a;s)\)
                \STATE \(\mathrm{trace}[i,h]\gets(h-b,a)\)
            \ENDIF
        \ENDFOR
    \ENDFOR
    \STATE \(D\gets D_{\mathrm{new}}\)
\ENDFOR
\STATE \(h\gets H\)
\FOR{\(i=M\) \TO \(1\)}
    \STATE \((h_{\mathrm{prev}},a_i)\gets\mathrm{trace}[i,h]\)
    \STATE \(h\gets h_{\mathrm{prev}}\)
\ENDFOR
\STATE \textbf{return} \(\pi=(a_1,\ldots,a_M)\)
\end{algorithmic}
\end{algorithm}

Algorithm~\ref{alg:budgeted-policy-optimization} shows the single-budget
solver. One run up to a maximum budget also produces the terminal states
\(D_M(h)\) for all \(h\leq H\), which can be backtracked to obtain multiple
frontier candidates. The two DP rows require \(O(H)\) working memory; retaining
all backpointers for direct reconstruction requires an additional \(O(MH)\)
trace, which can be avoided by recomputing rows during backtracking.

The proxy models and solver are used for candidate screening rather than as a
substitute for measurement. Reported operating points are exported using the
same canonical compression states used during calibration and are re-evaluated
with the real extended vLLM runtime and downstream tasks.

\FloatBarrier
\begin{figure*}[t]
    \centering
    \includegraphics[width=0.49\textwidth]{Figures/draft2/pareto_search_audit.pdf}
    \hfill
    \includegraphics[width=0.49\textwidth]{Figures/draft2/pareto_search_audit_display.pdf}
    \caption{Empirical audit of the measured DASO frontiers on
    Llama-3.1-8B-Instruct for prefill-only inference (left) and
    prefill--decode inference (right). Red squares denote uniform baselines;
    blue, green, and orange markers denote random, neighborhood, and
    leverage-guided challenger policies, respectively. The dashed curves are
    the frontiers after combining the challengers with the existing DASO
    points.}
    \label{fig:pareto-search-audit}
\end{figure*}
\FloatBarrier

\section{Empirical Pareto-Frontier Audit}
\label{app:pareto-audit}
This section is intended to address a supplementary question: how close are the policies found
by DASO to the true Pareto frontier? To probe this question, we attempt to
outperform the DASO points using several broader, independently constructed
search strategies, as described below.

\textbf{Audit setup.}
We examine whether the measured DASO frontier is merely favored by its
proxy-guided optimization path. For each of the prefill-only and
prefill--decode scenarios, we independently generate 30 label-blind challenger
policies: 15 random module-wise policies, 10 neighborhood mutations around
representative DASO points, and five leverage-guided mutations targeting
modules with promising predicted latency--quality ratios. Candidate generation
does not access the downstream labels. Every challenger is then exported and
fully evaluated using the same extended vLLM runtime and downstream protocols
as the reported policies. The comparison in
Fig.~\ref{fig:pareto-search-audit} therefore uses measured E2E speedup and full
ARC-Challenge or IWSLT quality, rather than proxy predictions.

\textbf{Results.}
Across the two scenarios, 18 of the 20 existing DASO frontier points remain
non-dominated after adding the 60 challengers. All seven prefill--decode points
are retained. In prefill-only inference, only two of 13 points are dominated,
and the largest observed improvement is \(0.0305\times\) in speedup and
\(0.005973\) in normalized ARC-Challenge accuracy. The challengers primarily
insert isolated bridge points or extend the low-quality tail, rather than
forming a systematically better frontier. These results indicate that the
DASO policies are already close to the empirical Pareto boundary discoverable
under this independent search budget and are robust to alternative search
paths. Although this finite audit cannot establish global optimality over the
combinatorial policy space, it provides empirical evidence that DASO approaches the underlying Pareto frontier.

\end{document}
